% Generated draft: requires mathematical and source review.
\documentclass[11pt]{article}
\usepackage[margin=25mm]{geometry}
\usepackage{fontspec}
\setmainfont{Latin Modern Roman}
\usepackage{amsmath,amssymb,mathtools,unicode-math}
\setmathfont{Latin Modern Math}
\usepackage{enumitem,needspace,xurl,hyperref}
\setlength{\parindent}{0pt}
\setlength{\parskip}{5pt}
\setlength{\emergencystretch}{3em}
\newcommand{\N}{\mathbb{N}}
\newcommand{\Z}{\mathbb{Z}}
\newcommand{\Q}{\mathbb{Q}}
\newcommand{\R}{\mathbb{R}}
\newcommand{\C}{\mathbb{C}}
\newcommand{\F}{\mathbb{F}}
\newcommand{\A}{\mathbb{A}}
\newcommand{\PP}{\mathbb{P}}
\newcommand{\E}{\mathbb{E}}
\newcommand{\eps}{\varepsilon}
\newcommand{\dd}{\,\mathrm{d}}
\begin{document}
\section*{UnsolvedMath ID: 30005771}
\textbf{Source problem code:} catalogue entry
\subsection{Definitions and mathematical statement}
For each positive integer \(n\), let \(S_n\) be the set of bijections from \(\{1,\ldots,n\}\) to itself, and let \(\pi_n\) be uniformly distributed on \(S_n\). The supplied fragment records the contextual fact
\[
\Pr\bigl(\pi_n(i)=j\bigr)=\frac{1}{n}
\qquad\text{for every }i,j\in\{1,\ldots,n\}.
\]
Thus these point probabilities are not \(o(1/n)\) as \(n\to\infty\): their ratio to \(1/n\) is identically \(1\).

\textbf{Unrecoverable intended statement.} The extracted passage begins with the isolated word ``conjecture'' and refers to an unspecified condition. It supplies neither that condition nor a complete conjecture or open question. The supplied statement-verification note reports that the passage comments on a proved point-probability criterion and does not state an additional conjecture or open question. The criterion, its hypotheses, its conclusion, and its asymptotic quantifiers are absent and cannot be reconstructed from this record. The displayed identity is contextual information, not a recovered open-problem statement.
\subsection{Short English statement}
A uniformly random permutation of size \(n\) sends each position to each value with probability exactly \(1/n\). The supplied passage comments on an omitted criterion; it does not provide a complete problem. The registry's supplied \texttt{open} label does not establish that this fragment states an open question.
\subsection{Sources}
\begin{itemize}
\item Local registry: \texttt{data/unsolvedmath/problems.json}, permanent UnsolvedMath ID \(30005771\), title ``Point-Probability Criterion Commentary,'' source problem code \texttt{catalogue entry}. The supplied record labels the status \texttt{open} and the statement status \texttt{unrecoverable}; no clean statement is supplied.
\item Supplied report citation: \emph{Mini-Workshop: Permutation Patterns} (2024), \emph{Oberwolfach Reports}, pp.\ 273--308, DOI \texttt{10.4171/owr/2024/6}. Source URL and supplied literature source labelled ``original OWR report'': \url{https://doi.org/10.4171/owr/2024/6}. No individual proposer is identified.
\item Supplied literature assessment, checked \texttt{2026-08-22}: the extracted text is only a fragment labelled ``conjecture'' and does not state a complete mathematical claim, so no authoritative resolution could be verified. This is supplied dated curation, not an independently verified literature assessment.
\item The supplied statement-verification note says that the official report confirms that the fragment comments on a proved point-probability criterion and states no additional conjecture or open question. It also says that no separate clean statement is promoted and that the original record remains publicly visible for provenance. This verification claim is supplied metadata and has not been independently checked here. The record's supplied update timestamp is \texttt{2026-09-16T12:06:53Z}.
\end{itemize}
\end{document}
